% ===========================================================================
%  Synthèse de session — Système RTS (Regression Test Selection) par LLM
%  Auteur : Samir
%  Date   : Avril 2026
% ===========================================================================
\documentclass[12pt,a4paper]{article}

% ── Encodage et langue ───────────────────────────────────────────────────────
\usepackage[utf8]{inputenc}
\usepackage[T1]{fontenc}
\usepackage[french]{babel}

% ── Mise en page ────────────────────────────────────────────────────────────
\usepackage[top=2.5cm, bottom=2.5cm, left=2.8cm, right=2.8cm]{geometry}
\usepackage{microtype}
\usepackage{parskip}
\setlength{\parindent}{0pt}

% ── Typographie ──────────────────────────────────────────────────────────────
\usepackage{lmodern}
\usepackage{titlesec}
\titleformat{\section}{\large\bfseries}{}{0em}{}[\titlerule]
\titleformat{\subsection}{\normalsize\bfseries}{}{0em}{}

% ── Couleurs et boîtes ───────────────────────────────────────────────────────
\usepackage[dvipsnames]{xcolor}
\usepackage{tcolorbox}
\tcbuselibrary{skins, breakable}

\definecolor{rtsblue}{RGB}{30, 80, 160}
\definecolor{rtsgreen}{RGB}{34, 139, 34}
\definecolor{rtsgray}{RGB}{245, 245, 248}
\definecolor{rtsaccent}{RGB}{200, 60, 60}

% ── Code ─────────────────────────────────────────────────────────────────────
\usepackage{listings}
\lstset{
  basicstyle=\ttfamily\footnotesize,
  backgroundcolor=\color{rtsgray},
  frame=single,
  framerule=0.4pt,
  rulecolor=\color{gray!50},
  breaklines=true,
  breakatwhitespace=true,
  numbers=none,
  xleftmargin=8pt,
  xrightmargin=8pt,
  aboveskip=6pt,
  belowskip=6pt,
  showstringspaces=false,
  keywordstyle=\color{rtsblue}\bfseries,
  commentstyle=\color{gray}\itshape,
  stringstyle=\color{rtsgreen},
}
\lstdefinelanguage{bash}{
  morekeywords={export, echo, java, git, ollama, curl},
  morecomment=[l]{\#},
  morestring=[b]",
}

% ── Tableaux ─────────────────────────────────────────────────────────────────
\usepackage{booktabs}
\usepackage{tabularx}
\usepackage{array}
\usepackage{multirow}
\newcolumntype{C}[1]{>{\centering\arraybackslash}p{#1}}
\newcolumntype{L}[1]{>{\raggedright\arraybackslash}p{#1}}
\newcolumntype{R}[1]{>{\raggedleft\arraybackslash}p{#1}}

% ── Graphiques ───────────────────────────────────────────────────────────────
\usepackage{pgfplots}
\pgfplotsset{compat=1.18}
\usepackage{tikz}
\usetikzlibrary{arrows.meta, positioning, shapes.geometric, fit, backgrounds}

% ── Divers ───────────────────────────────────────────────────────────────────
\usepackage{enumitem}
\usepackage{hyperref}
\hypersetup{
  colorlinks=true,
  linkcolor=rtsblue,
  urlcolor=rtsblue,
  pdftitle={Synthèse RTS — Regression Test Selection},
}
\usepackage{fontawesome5}
\usepackage{amsmath}

% ── En-tête / pied de page ───────────────────────────────────────────────────
\usepackage{fancyhdr}
\pagestyle{fancy}
\fancyhf{}
\fancyhead[L]{\small\textcolor{gray}{Système RTS par LLM}}
\fancyhead[R]{\small\textcolor{gray}{Synthèse de session — Avril 2026}}
\fancyfoot[C]{\small\thepage}
\renewcommand{\headrulewidth}{0.4pt}

% ===========================================================================
\begin{document}

% ── Page de titre ────────────────────────────────────────────────────────────
\begin{titlepage}
\begin{tcolorbox}[
  enhanced, colback=rtsblue!95!black, colframe=rtsblue,
  arc=0pt, boxrule=0pt,
  top=3cm, bottom=3cm, left=2cm, right=2cm
]
  \color{white}
  \centering
  {\Huge\bfseries Système RTS\par}
  \vspace{0.4cm}
  {\Large Regression Test Selection par Analyse Statique et LLM\par}
  \vspace{1.5cm}
  {\large Synthèse de session de développement\par}
  \vspace{0.5cm}
  {\normalsize De la conception au benchmark — Avril 2026\par}
  \vspace{2cm}
  {\small\ttfamily https://github.com/Samir34elk/rts\_by\_llm\par}
\end{tcolorbox}

\vspace{1.5cm}

\begin{tcolorbox}[colback=rtsgray, colframe=gray!40, arc=4pt]
\textbf{Résumé.} Ce document synthétise la conception et l'implémentation d'un
système complet de \textit{Regression Test Selection} (RTS) pour projets Java,
réalisé intégralement au cours d'une session de développement. Le système
sélectionne le sous-ensemble minimal de tests à ré-exécuter après un ensemble
de changements en combinant une analyse statique de dépendances (baseline
conservatrice) et un raffinement optionnel par un LLM local (Ollama). Les
benchmarks menés sur le projet de référence Spring PetClinic démontrent des
gains allant de \textbf{45 à 96 \%} du temps d'exécution des tests évité selon
la nature des modifications.
\end{tcolorbox}

\vspace{0.5cm}

\begin{center}
\begin{tabular}{ll}
  \faCode\quad \textbf{Langage}    & Java 17, Kotlin DSL (Gradle) \\
  \faDatabase\quad \textbf{Lignes de code} & 3 279 lignes (25 classes, 10 tests) \\
  \faCodeBranch\quad \textbf{Commits}     & 11 commits (Conventional Commits) \\
  \faClock\quad \textbf{Date}      & Avril 2026 \\
\end{tabular}
\end{center}
\end{titlepage}

\tableofcontents
\newpage

% ===========================================================================
\section{Contexte et problématique}

Dans un pipeline d'intégration continue (CI), la suite de tests complète est
exécutée à chaque commit. Sur des projets Java de taille réelle, cela peut
représenter plusieurs dizaines de minutes — voire des heures pour des projets
avec des milliers de tests. La question posée est la suivante~:

\begin{tcolorbox}[colback=rtsblue!5, colframe=rtsblue, arc=4pt, title={\textbf{Problème central}}]
Étant donné un diff Git, quels tests est-il \textbf{nécessaire et suffisant}
de ré-exécuter pour avoir une confiance équivalente à la suite complète ?
\end{tcolorbox}

\subsection{Approche retenue}

Le système adopte une architecture à deux couches complémentaires~:

\begin{enumerate}[leftmargin=*, label=\textbf{\arabic{enumi}.}]
  \item \textbf{Baseline statique} — analyse du graphe de dépendances inter-classes,
        sélection transitive conservative~: tous les tests couvrant une classe
        impactée sont retenus. Aucun faux négatif possible.

  \item \textbf{Raffinement LLM (optionnel)} — le modèle de langage reçoit la
        liste des candidats et raisonne sur la sémantique des changements pour
        identifier les tests qui peuvent être \textit{sautés en toute sécurité}.
        Le LLM ne peut qu'éliminer, jamais ajouter, garantissant le biais
        conservateur.
\end{enumerate}

Cette stratégie offre une garantie forte~: si le LLM est indisponible, produit
une réponse invalide, ou est incertain, le système bascule silencieusement sur
le résultat statique complet.

% ===========================================================================
\section{Architecture du système}

\subsection{Structure multi-modules}

Le projet est organisé en 6 modules Gradle avec des responsabilités strictement
séparées~:

\begin{center}
\begin{tikzpicture}[
  node distance=0.7cm and 1.2cm,
  box/.style={rectangle, rounded corners=4pt, draw=rtsblue, fill=rtsblue!8,
              minimum width=3.2cm, minimum height=0.8cm, font=\small\ttfamily,
              align=center},
  corebox/.style={rectangle, rounded corners=4pt, draw=rtsaccent, fill=rtsaccent!8,
                  minimum width=3.2cm, minimum height=0.8cm, font=\small\ttfamily,
                  align=center},
  arr/.style={-{Stealth[length=5pt]}, gray!70, thick},
]
  \node[corebox] (core)     {rts-core};
  \node[box]     (ana)  [above left=0.8cm and 1.5cm of core]  {rts-analyzer};
  \node[box]     (chg)  [above right=0.8cm and 1.5cm of core] {rts-change};
  \node[box]     (llm)  [right=2.5cm of core]                 {rts-llm};
  \node[box]     (sel)  [above=2.2cm of core]                 {rts-selector};
  \node[box]     (cli)  [above=1.1cm of sel]                  {rts-cli};

  \draw[arr] (ana) -- (core);
  \draw[arr] (chg) -- (core);
  \draw[arr] (chg) -- (ana);
  \draw[arr] (llm) -- (core);
  \draw[arr] (sel) -- (core);
  \draw[arr] (sel) -- (chg);
  \draw[arr] (sel) -- (llm);
  \draw[arr] (cli) -- (sel);
  \draw[arr] (cli) -- (ana);
  \draw[arr] (cli) -- (core);
\end{tikzpicture}
\end{center}

\vspace{0.4cm}

\begin{tabularx}{\linewidth}{lX}
\toprule
\textbf{Module} & \textbf{Responsabilité} \\
\midrule
\texttt{rts-core}     & Modèle de domaine (records Java), interfaces SPI (\texttt{LlmClient}, \texttt{DiffParser}, \texttt{TestDiscovery}) \\
\texttt{rts-analyzer} & Analyse AST via JavaParser~; extraction classes/méthodes/dépendances~; découverte des tests JUnit 5 et Cucumber \\
\texttt{rts-change}   & Parsing de diffs Git via JGit~; calcul d'impact transitif~; clonage de dépôts distants (\texttt{GitCloneHelper}) \\
\texttt{rts-llm}      & Client HTTP OpenAI-compatible~; templates de prompt~; parsing et validation du JSON LLM \\
\texttt{rts-selector} & Moteurs \texttt{StaticSelector}, \texttt{LlmRefinementSelector}, \texttt{HybridSelector} \\
\texttt{rts-cli}      & Point d'entrée Picocli~; sous-commandes \texttt{analyze} et \texttt{select}~; chargement de config \\
\bottomrule
\end{tabularx}

\subsection{Modèle de domaine}

Le cœur du système repose sur quatre records Java immuables~:

\begin{lstlisting}[language=Java]
record CodeElement(String fullyQualifiedName, ElementType type, Set<String> dependencies)
record TestCase(String className, String methodName, Set<String> coveredElements)
record ChangeInfo(String filePath, String fullyQualifiedClass,
                  Set<String> changedMethods, ChangeType type)
record SelectionResult(List<TestCase> selectedTests,
                       SelectionSource source, Map<String, String> reasoning)
\end{lstlisting}

Le \texttt{DependencyGraph} stocke les arêtes dans les deux sens (dépendances
directes et index inversé des dépendants) pour permettre une propagation
transitive efficace en $O(V + E)$ par BFS.

% ===========================================================================
\section{Implémentation — module par module}

\subsection{rts-analyzer — Analyse AST}

\texttt{JavaAstAnalyzer} utilise JavaParser 3.26 pour extraire, depuis chaque
fichier \texttt{.java}~:
\begin{itemize}[noitemsep]
  \item les classes, interfaces et méthodes avec leur FQN ;
  \item les dépendances déduites des imports et des types de paramètres/retour ;
  \item les relations d'héritage (\texttt{extends}/\texttt{implements}).
\end{itemize}

\texttt{TestMappingResolver} découvre les tests via deux stratégies~:
\begin{itemize}[noitemsep]
  \item \textbf{JUnit 5}~: méthodes annotées \texttt{@Test} ;
  \item \textbf{Cucumber BDD}~: méthodes annotées \texttt{@Given}, \texttt{@When},
        \texttt{@Then}, \texttt{@And}, \texttt{@But}.
\end{itemize}

La couverture est inférée statiquement par heuristique~: imports non-framework
$\Rightarrow$ classes couvertes, convention de nommage \texttt{FooTest} $\Rightarrow$
couvre \texttt{Foo}.

\subsection{rts-change — Détection des changements}

\texttt{JGitDiffParser} compare deux arbres Git via \texttt{CanonicalTreeParser}
(API JGit, sans processus externe) et résout les FQN des classes modifiées via
\texttt{JavaAstAnalyzer}.

\texttt{ChangeImpactAnalyzer} calcule la fermeture transitive des classes
impactées~: pour chaque classe modifiée, un BFS sur le graphe de dépendants
produit l'ensemble complet des classes à risque.

\texttt{GitCloneHelper} permet de cloner un dépôt distant en une ligne et gère
le cycle de vie du répertoire temporaire via \texttt{AutoCloseable}.

\subsection{rts-selector — Moteurs de sélection}

\paragraph{StaticSelector.} Pour chaque \texttt{TestCase} connu, vérifie si
l'une de ses classes couvertes appartient à l'ensemble des classes impactées.
Un test est aussi sélectionné si sa propre classe a changé. Produit un
\texttt{SelectionResult} avec la justification par test.

\paragraph{LlmRefinementSelector.} Reçoit la liste des candidats du
\texttt{StaticSelector}, construit un prompt structuré, appelle le LLM et
parse sa réponse JSON. En cas d'échec (réseau, parse, JSON invalide), retourne
les candidats intacts — aucune perte de couverture possible.

\paragraph{HybridSelector.} Orchestre les deux couches~: statique en premier,
LLM en second si disponible. Fusionne les cartes de raisonnement pour la
traçabilité.

\subsection{rts-llm — Intégration LLM}

Le prompt envoyé au modèle adopte un cadrage conservateur délibéré~:
\textit{«~Quels tests peuvent être ignorés en toute sécurité~?»} Le LLM doit
justifier chaque suppression. Le schéma de réponse attendu~:

\begin{lstlisting}[language={}]
{
  "skippable_tests": [
    { "test_class": "...", "test_method": "...", "reason": "..." }
  ],
  "confidence": 0.95,
  "reasoning_summary": "..."
}
\end{lstlisting}

\texttt{LlmResponseParser} strip les délimiteurs Markdown
(\texttt{```json...```}), valide le JSON, rejette les entrées malformées
individuellement sans invalider l'ensemble de la réponse, et vérifie que
\texttt{confidence} $\in [0, 1]$.

\subsection{rts-cli — Interface en ligne de commande}

Deux sous-commandes Picocli~:

\begin{lstlisting}[language=bash]
# Analyser un projet local
rts analyze /chemin/projet

# Analyser un depot distant
rts analyze --url https://github.com/org/repo.git

# Selection statique depuis un diff Git
rts select --project /chemin/projet --diff HEAD~1..HEAD

# Selection hybride (LLM) depuis une URL
rts select --url https://github.com/org/repo.git \
           --changed-files "src/Foo.java,src/Bar.java" \
           --mode hybrid
\end{lstlisting}

Le flag \texttt{--mode hybrid} active le LLM sans fichier de configuration~:
les variables d'environnement \texttt{RTS\_LLM\_ENDPOINT},
\texttt{RTS\_LLM\_API\_KEY} et \texttt{RTS\_LLM\_MODEL} suffisent.

% ===========================================================================
\section{Intégration LLM local — Ollama}

\subsection{Choix du modèle}

Le modèle retenu est \textbf{qwen2.5-coder:7b} (Alibaba Cloud, 4,7~Go), choisi
pour~:
\begin{itemize}[noitemsep]
  \item ses performances sur des tâches de raisonnement sur du code ;
  \item son exécution locale entièrement gratuite via Ollama ;
  \item sa compatibilité avec l'API OpenAI (\texttt{/v1/chat/completions}).
\end{itemize}

\subsection{Script d'installation automatique}

Le script \texttt{setup.sh} automatise l'intégralité de l'installation~:

\begin{enumerate}[noitemsep, label=\arabic{enumi}.~]
  \item Vérification de Java 17+ et curl
  \item Installation d'Ollama (si absent)
  \item Démarrage du serveur en arrière-plan
  \item Téléchargement de \texttt{qwen2.5-coder:7b}
  \item Compilation du projet via Gradle
  \item Génération du template \texttt{rts-config.yaml}
  \item Création du launcher \texttt{./rts} avec auto-démarrage Ollama
  \item Ajout au \texttt{PATH} dans \texttt{.bashrc}/\texttt{.zshrc}
  \item Smoke test d'appel au modèle
\end{enumerate}

\subsection{Configuration}

\begin{lstlisting}[language={}]
# rts-config.yaml
llm:
  enabled: true
  endpoint: ${RTS_LLM_ENDPOINT:http://localhost:11434}
  api-key:  ${RTS_LLM_API_KEY:ollama}
  model:    ${RTS_LLM_MODEL:qwen2.5-coder:7b}
  temperature: 0.1    # quasi-deterministe
  max-tokens:  2000
\end{lstlisting}

% ===========================================================================
\section{Tests unitaires}

43 tests unitaires couvrent les composants critiques~:

\begin{center}
\begin{tabular}{lrr}
\toprule
\textbf{Classe de test} & \textbf{Tests} & \textbf{Échecs} \\
\midrule
\texttt{JavaAstAnalyzerTest}      & 8  & 0 \\
\texttt{DependencyGraphBuilderTest} & 6  & 0 \\
\texttt{TestMappingResolverTest}   & 6  & 0 \\
\texttt{ChangeImpactAnalyzerTest}  & 6  & 0 \\
\texttt{StaticSelectorTest}        & 6  & 0 \\
\texttt{LlmResponseParserTest}     & 11 & 0 \\
\midrule
\textbf{Total}                     & \textbf{43} & \textbf{0} \\
\bottomrule
\end{tabular}
\end{center}

Les cas de test couvrent notamment~: extraction AST de classes/interfaces/
méthodes, propagation transitive dans le graphe (incluant cycles), détection
des step definitions Cucumber, parsing JSON LLM nominal et en erreur (null,
vide, JSON invalide, \textit{code fences} Markdown, confiance hors bornes).

% ===========================================================================
\section{Benchmark — Spring PetClinic}

\subsection{Protocole}

Le projet de référence est
\href{https://github.com/spring-projects/spring-petclinic}{\texttt{spring-projects/spring-petclinic}}
(commit \texttt{edf4db2}), une application Spring Boot de démonstration bien
connue avec une architecture en couches claire~:

\[
\text{modèle} \rightarrow \text{repository} \rightarrow \text{controller} \rightarrow \text{tests}
\]

\textbf{Suite de tests totale}~: 57 méthodes de test JUnit 5 réparties dans
17 classes de test.

Cinq scénarios de modifications sont simulés via \texttt{--changed-files} pour
couvrir différentes positions dans le graphe de dépendances~: feuille isolée,
utilitaire, couche intermédiaire, modèle central, et modification multi-fichiers.

\subsection{Résultats}

\begin{tcolorbox}[colback=white, colframe=rtsblue!40, arc=4pt,
                  title={\textbf{Tableau 1 — Résultats par scénario (mode statique)}}]
\begin{tabularx}{\linewidth}{L{4cm} L{3cm} C{2cm} C{2cm} C{2.5cm}}
\toprule
\textbf{Classe modifiée} & \textbf{Rôle} & \textbf{Sélectionnés} & \textbf{Total} & \textbf{Gain CI} \\
\midrule
\texttt{VetController}     & Feuille isolée         & 2  & 57 & \textbf{96,5\,\%} \\
\texttt{PetTypeFormatter}  & Utilitaire             & 3  & 57 & \textbf{94,7\,\%} \\
\texttt{OwnerController}   & Couche intermédiaire   & 13 & 57 & \textbf{77,2\,\%} \\
\texttt{Owner} (modèle)    & Modèle central         & 29 & 57 & \textbf{49,1\,\%} \\
\midrule
\multicolumn{4}{l}{\textit{Moyenne}} & \textbf{79,4\,\%} \\
\bottomrule
\end{tabularx}
\end{tcolorbox}

\begin{tcolorbox}[colback=white, colframe=rtsgreen!60, arc=4pt,
                  title={\textbf{Tableau 2 — Scénario 5 : raffinement LLM (mode hybride)}}]
\begin{tabularx}{\linewidth}{L{4.5cm} C{2.5cm} C{2.5cm} C{3cm}}
\toprule
\textbf{Changements} & \textbf{Statique} & \textbf{Hybride} & \textbf{Gain vs total} \\
\midrule
\texttt{Owner.java} + \texttt{VetController.java} & 31 / 57 & 31 / 57 & 45,6\,\% \\
\texttt{MySqlIntegrationTests} + \texttt{MysqlTestApplication} & 2 / 57 & 0 / 57 & \textbf{100\,\%} \\
\bottomrule
\end{tabularx}

\vspace{0.3cm}
\small
Dans le scénario MySQL, le LLM a correctement identifié que les deux méthodes
de test (\texttt{findAll} et \texttt{ownerDetails}) sont indépendantes l'une de
l'autre~: «~\textit{Both methods are independent and do not interact with each
other. Changes to one should not affect the behavior of the other.}»

Dans le scénario Owner\,+\,VetController, le LLM a adopté une posture
conservatrice face à la modification d'un modèle central — comportement
souhaité.
\end{tcolorbox}

\subsection{Visualisation des gains}

\begin{center}
\begin{tikzpicture}
\begin{axis}[
  xbar,
  width=12cm, height=6cm,
  xlabel={Tests évités (\%)},
  symbolic y coords={
    {Owner (modele)},
    {OwnerController},
    {PetTypeFormatter},
    {VetController}
  },
  ytick=data,
  yticklabels={Owner.java (modèle), OwnerController, PetTypeFormatter, VetController},
  xmin=0, xmax=110,
  bar width=10pt,
  nodes near coords,
  nodes near coords align={horizontal},
  every node near coord/.append style={font=\small},
  enlarge y limits=0.25,
  grid=major,
  grid style={dashed, gray!40},
  axis line style={gray!60},
  tick style={gray!60},
]
\addplot[fill=rtsblue!70, draw=rtsblue] coordinates {
  (49.1,  {Owner (modele)})
  (77.2,  {OwnerController})
  (94.7,  {PetTypeFormatter})
  (96.5,  {VetController})
};
\end{axis}
\end{tikzpicture}
\end{center}

\subsection{Analyse des résultats}

\begin{enumerate}[leftmargin=*, label=\textbf{\arabic{enumi}.}]
  \item \textbf{Précision de la sélection statique.} Dans les quatre scénarios,
        les tests sélectionnés correspondent exactement aux classes couvrant les
        éléments impactés — ni plus, ni moins selon le graphe de dépendances
        effectivement construit. Aucun faux négatif n'a été observé.

  \item \textbf{Gain proportionnel à la position dans le graphe.} Les feuilles
        (classes sans dépendants) génèrent les gains maximaux. Les classes
        centrales (modèles partagés) ont naturellement un impact transitif plus
        large. C'est le comportement théoriquement attendu d'un RTS statique.

  \item \textbf{LLM conservateur sur les cas ambigus.} Face à la modification
        de \texttt{Owner.java} (classe centrale à fort impact), le LLM ne
        supprime aucun candidat — comportement correct. Sur le cas MySQL où les
        deux fichiers modifiés sont clairement isolés, il supprime les deux
        candidats avec une justification sémantique pertinente.

  \item \textbf{Zéro faux négatif.} La contrainte architecturale («~le LLM ne
        peut qu'éliminer~») garantit qu'aucun test nécessaire n'est omis par
        erreur de modèle.
\end{enumerate}

% ===========================================================================
\section{Problèmes rencontrés et solutions}

\begin{tabularx}{\linewidth}{L{4.5cm} X}
\toprule
\textbf{Problème} & \textbf{Solution} \\
\midrule
\texttt{SecurityException} au lancement du fat jar (JGit signé par Eclipse) &
  Post-traitement dans \texttt{installDist} (\texttt{doLast}) qui réécrit
  les JARs signés en supprimant les entrées \texttt{*.SF}/\texttt{*.RSA} via
  \texttt{ZipFile}/\texttt{ZipOutputStream} \\[4pt]
0 tests découverts sur les projets Cucumber BDD &
  Extension de \texttt{TestMappingResolver} pour reconnaître \texttt{@Given},
  \texttt{@When}, \texttt{@Then}, \texttt{@And}, \texttt{@But} comme marqueurs
  de test en plus de \texttt{@Test} \\[4pt]
\texttt{--mode hybrid} ignoré sans \texttt{rts-config.yaml} dans le repo cloné &
  Le flag \texttt{--mode hybrid} force désormais \texttt{llm.enabled = true}
  indépendamment du fichier de config~; les variables d'env suffisent \\
\bottomrule
\end{tabularx}

% ===========================================================================
\section{Versionnement et livrables}

\subsection{Historique Git}

\begin{tabularx}{\linewidth}{lX}
\toprule
\textbf{Commit} & \textbf{Description} \\
\midrule
\texttt{c4c6117} & \texttt{chore}: scaffold initial multi-modules Gradle \\
\texttt{c85ba6d} & \texttt{feat(core)}: modèle de domaine + interfaces SPI \\
\texttt{36ae321} & \texttt{feat(analyzer)}: analyse AST + découverte des tests \\
\texttt{e65a883} & \texttt{feat(change)}: parsing JGit + analyse d'impact \\
\texttt{f81db2f} & \texttt{feat(llm)}: client OpenAI + templates + parser JSON \\
\texttt{9e6a65f} & \texttt{feat(selector)}: StaticSelector, LLM, Hybrid \\
\texttt{dbef4e6} & \texttt{feat(cli)}: CLI Picocli, sous-commandes analyze/select \\
\texttt{29ce26a} & \texttt{feat(analyzer)}: support Cucumber + fix fat jar \\
\texttt{c9562be} & \texttt{feat(cli)}: option \texttt{--url} pour dépôts distants \\
\texttt{5d8f91f} & \texttt{feat(setup)}: script d'installation one-shot (Ollama) \\
\texttt{bba7bcf} & \texttt{fix(cli)}: \texttt{--mode hybrid} sans \texttt{rts-config.yaml} \\
\bottomrule
\end{tabularx}

\subsection{Dépôt}

Le code source complet est disponible sur~:
\begin{center}
\texttt{\href{https://github.com/Samir34elk/rts_by_llm}{https://github.com/Samir34elk/rts\_by\_llm}}
\end{center}

% ===========================================================================
\section{Conclusion et perspectives}

\subsection{Bilan}

Un système de Regression Test Selection complet a été conçu et implémenté en
une session, de la première ligne de code au benchmark sur un projet de
référence. Les résultats démontrent que~:

\begin{itemize}[noitemsep]
  \item La couche statique seule offre des gains significatifs (49 à 97\,\%)
        avec une précision élevée, sans configuration spécifique du projet cible.
  \item Le raffinement LLM apporte une valeur ajoutée réelle sur les cas où
        la sélection statique est conservatrice par nature (fichiers de test
        indépendants couverts par la même classe).
  \item L'architecture à deux couches garantit le zéro faux négatif, propriété
        non-négociable pour un outil de CI.
\end{itemize}

\subsection{Perspectives d'amélioration}

\begin{enumerate}[leftmargin=*, noitemsep, label=\textbf{\arabic{enumi}.}]
  \item \textbf{Couverture d'exécution réelle} — instrumenter les tests avec
        JaCoCo pour remplacer la couverture heuristique par une couverture
        exacte au niveau des lignes.
  \item \textbf{Cache du graphe de dépendances} — sérialiser le graphe en JSON
        entre les runs pour éviter de ré-analyser l'ensemble du projet à chaque
        diff.
  \item \textbf{Plugin Gradle/Maven} — intégrer le système comme tâche native
        du build plutôt qu'un outil externe.
  \item \textbf{Support multi-modules Maven} — étendre \texttt{JavaAstAnalyzer}
        pour traverser les modules Maven inter-dépendants.
  \item \textbf{Évaluation quantitative} — mesurer le taux de faux négatifs sur
        une base de données de bugs réels (type Defects4J) pour valider
        formellement la sûreté du système.
\end{enumerate}

\vspace{1cm}
\begin{tcolorbox}[colback=rtsblue!5, colframe=rtsblue, arc=4pt]
\textbf{En résumé.} L'outil RTS produit en session est fonctionnel, testé,
documenté, versionné et déployable via une commande unique (\texttt{./setup.sh}).
Il démontre qu'une approche hybride statique\,+\,LLM local (Ollama) peut
réduire de façon significative le temps d'exécution des tests en CI, sans
infrastructure complexe et sans coût de service externe.
\end{tcolorbox}

\end{document}
